% !TeX root = ../main.tex

\chapter{System Overview and Unified Development Framework}

\section{Hexapole Actuation Platform}

This dissertation considers a Hall-sensor-based hexapole electromagnetic actuator for precision manipulation in aqueous solution. The platform uses six independently driven poles to generate magnetic flux and magnetic-field gradients in a three-dimensional workspace. Because force objectives are three-dimensional while actuation is six-dimensional, the system is over-actuated and requires explicit allocation and control logic.

From an engineering perspective, the platform combines non-contact operation, strong force authority, and microscope compatibility. These properties make it suitable for high-speed untethered manipulation when flux generation and multi-axis interaction are properly handled.

\section{End-to-End Control Chain}

The physical and computational chain considered in this work starts from actuation commands, proceeds through magnetic-flux generation, and ends at force output on the probe. Hall-sensor measurements provide the core feedback channel for closed-loop flux regulation, and the force-generation stage maps control outputs to task-level force behavior.

For dissertation-level reasoning, this chain is treated as one connected pipeline rather than isolated blocks. Upstream uncertainty and design decisions can propagate to downstream force quality, so each stage must be interpreted in relation to the full chain.

\section{Development-Process Structuring Blueprint (Pillar A)}

The first equal-weight pillar is process structuring. Chapter 2 defines a reusable blueprint for developing hexapole-form tweezers:

1. establish platform and measurement context;
2. extract control-use model information from calibration data;
3. design and validate flux-control structure;
4. evaluate force-generation quality under design choices;
5. close the chain with cross-layer interpretation.

This blueprint makes stage boundaries explicit and keeps design decisions traceable from observation to final force behavior.

\section{Workflow Automation Interfaces (Pillar B)}

The second equal-weight pillar is workflow automation. Chapter 2 defines where automation is integrated into the blueprint:

1. identification-stage repeatable processing and fitting flow;
2. control-stage simulation-and-tuning loop after identification;
3. stage-to-stage data handoff for force-generation validation.

These interfaces are defined to reduce development burden and improve reproducibility when deploying the framework on additional hexapole-form platforms.

\section{Chapter Interface to Chapters 3--5}

Chapter 2 provides the common language and unified framework for the remaining technical chapters. Chapter 3 develops the identification thread, Chapter 4 develops the flux-control thread, and Chapter 5 develops the force-generation thread and closes the cross-layer evidence chain.
